\section{Relation to the entropy--tail selection law}\label{sec:comparison}
This section is deterministic: it compares the closed functional
$\Q(B)$ with the selection functional of the companion general-law
manuscript~\cite{gl}. It is not used in the proof of
Theorem~\ref{thm:direct-main}.

\subsection{The selection functional}
Let $\mu$ be any probability law on $(0,1/2]$, not necessarily finite.
For such $\mu$, define $R_{\mu,\rho}(B)$ by \eqref{eq:minimax-risk},
with the model of Section~\ref{sec:model} unchanged except that
$D_k\sim\mu$, $q$ ranges over measurable functions with $0\le q\le1$,
and $\theta(q)=\int q\,d\mu$. Use the bands $I_j=(b_j/2,b_j]$, $b_j=2^{-j-1}$, of
\eqref{eq:core-entropy} and the functional $\mathcal G$ of
Section~\ref{sec:entropy-preliminaries}. For a finite measure
$\sigma\le\mu$ put $a_j(\sigma)=\int_{I_j}d^{-2}\,\sigma(dd)$ and
$A_\sigma=\sum_ja_j(\sigma)$, and define
$\Gamma_\sigma=\sup_J\mathcal G(a_0(\sigma),\ldots,a_J(\sigma))$,
with values in $[0,\infty]$. When only finitely many $a_j(\sigma)$ are
positive this is $\mathcal G(a(\sigma))$; for $\sigma=\kappa$ and a finite
law it is the core entropy of \eqref{eq:core-entropy}. Since $\mathcal G$
is coordinatewise increasing, $\sigma\le\sigma'$ implies
$\Gamma_\sigma\le\Gamma_{\sigma'}$, and $\Gamma_\sigma\ge A_\sigma$.
Write $H_\sigma=\int d^{-1}\,\sigma(dd)$ and, with $h=B^{-1/2}$,
\begin{equation}\label{eq:selection-functional}
 \mathcal K(B)=\min\left\{\rho,\ \inf_{0\le\sigma\le\mu}
 \left[h\sqrt{\rho\Gamma_\sigma}
       +\int\min\{1,h/d\}\,(\mu-\sigma)(dd)\right]\right\}.
\end{equation}
This is the functional $\mathcal K_{\mu,\rho}(B)$ of~\cite{gl}. That
manuscript states and proves $R_{\mu,\rho}(B)\asymp\mathcal K_{\mu,\rho}(B)$
for every probability law on $(0,1/2]$, every $0<\rho\le1/4$ and every
$B>0$, with explicit constants $10^{-17}$ and $33000$. Its proof has not
been externally refereed or formally verified, and it is not assumed
anywhere in this manuscript.

For a general law, define the quantile split as in
Section~\ref{sec:model}. Let $t=\min\{s:\mu((0,s])\ge\rho\}$, which
exists and is positive. Let $\nu$ be $\mu$ on $(0,t)$ plus the fraction
of the atom at $t$, if any, that gives $\nu$ total mass $\rho$, and let
$\kappa=\mu-\nu$. Then $\nu$ is supported on $(0,t]$, $\kappa$ is
supported on $[t,1/2]$, and only finitely many bands meet $[t,1/2]$.
Define $T_h=\int\min\{1,h/d\}\,\nu(dd)$ and
$\Q(B)=\min\{\rho,T_h+h\sqrt{\rho\Gamma_\kappa}\}$. For a finite law these
are the quantities of \eqref{eq:quantile-split}
and \eqref{eq:closed-functional}.

\subsection{Scale-separated cost}
Besides \eqref{eq:binary-mixture}, the comparison uses one more entropy
bound.

\begin{lemma}[Scale-separated cost]\label{lem:scale-separated}
If $t>0$ and $\xi$ is a finite measure on $[t,1/2]$, then
$\Gamma_\xi\le4H_\xi/t$.
\end{lemma}

\begin{proof}
Let $j_0$ be the band containing $t$, so $t\le b_{j_0}$. A band
meeting $[t,1/2]$ has index $j_0-r$ with $r\ge0$ and consists of gaps
in $(b_{j_0}2^{r-1},b_{j_0}2^r]$. There every gap of $\xi$ satisfies
$d\ge t$ and $d>b_{j_0}2^{r-1}\ge t2^{r-1}$, so $d^{-1}\le2^{1-r}/t$.
Writing $a_r$ and $H_r$ for the inverse-square and inverse costs of
$\xi$ in that band, $a_r\le2^{1-r}H_r/t$. Comparing the normalized
costs with the probability vector $\pi_r=2^{-(r+1)}$ gives
$\Ent(a/A)\le\sum_r(a_r/A)\log(1/\pi_r)$, whence
\[
 \Gamma_\xi\le\sum_{r\ge0}a_r[1+(r+1)\log2]
 \le\frac1t\sum_{r\ge0}2^{1-r}[1+(r+1)\log2]\,H_r.
\]
The coefficient $2^{1-r}[1+(r+1)\log2]$ is largest at $r=0$, where it
equals $2+2\log2<4$.
\end{proof}

\subsection{The comparison}
\begin{proposition}[Quantile versus selection]\label{prop:quantile-selection}
For every probability law $\mu$ on $(0,1/2]$, every $0<\rho\le1/4$ and
every $B>0$,
\[
 \mathcal K(B)\le\Q(B)\le(1+2\sqrt2)\,\mathcal K(B)<4\,\mathcal K(B).
\]
If $h\ge t$, then $\mathcal K(B)=\Q(B)=\rho$.
\end{proposition}

\begin{proof}
The measure $\sigma=\kappa$ is admissible in
\eqref{eq:selection-functional}, and $\mu-\kappa=\nu$. Its objective is
$h\sqrt{\rho\Gamma_\kappa}+T_h$, so $\mathcal K(B)\le\Q(B)$.

For the upper comparison fix $0\le\sigma\le\mu$ with finite objective,
let $f=d\sigma/d\mu\in[0,1]$, and write
\[
 x=h\sqrt{\rho\Gamma_\sigma},\qquad
 y=\int\min\{1,h/d\}\,(\mu-\sigma)(dd),\qquad S=x+y.
\]
Since $\nu\le\mu$, we have $(1-f)\nu\le\mu-\sigma$ and $f\nu\le\sigma$,
so $\int f\,d\nu\le\rho$ and $\int d^{-2}f\,\nu(dd)\le A_\sigma\le\Gamma_\sigma$.
With Cauchy--Schwarz,
\[
 T_h\le y+h\int d^{-1}f\,\nu(dd)
 \le y+h\sqrt{\rho\int d^{-2}f\,\nu(dd)}\le y+x=S.
\]
As $T_h\le\rho$ always, $T_h\le\mathcal K(B)$. If $h\ge t$, every gap of
$\nu$ has $\min\{1,h/d\}=1$, so $T_h=\rho$ and therefore
$\mathcal K(B)=\Q(B)=\rho$.

Now let $h<t$, and split $\kappa=f\kappa+(1-f)\kappa$. Both parts live
on the finitely many bands meeting $[t,1/2]$. By \eqref{eq:binary-mixture},
monotonicity and Lemma~\ref{lem:scale-separated},
\[
 \Gamma_\kappa\le2\Gamma_{f\kappa}+2\Gamma_{(1-f)\kappa}
 \le2\Gamma_\sigma+\frac8tH_{(1-f)\kappa}.
\]
Every gap of $\nu$ is at most $t$, so $\min\{1,h/d\}\ge h/t$ there and
$h\rho/t\le T_h$. Every gap of $\kappa$ is at least $t>h$, so
$hH_{(1-f)\kappa}=\int(h/d)\,(1-f)\kappa(dd)\le y$. Hence
\[
 h\sqrt{\rho\Gamma_\kappa}\le\sqrt2\,x
 +\sqrt{8\,(h\rho/t)\,(hH_{(1-f)\kappa})}\le\sqrt2\,x+\sqrt{8T_hy},
\]
and, with $T_h\le S$,
\[
 \Q(B)\le T_h+\sqrt2\,x+\sqrt{8T_hy}\le S+\sqrt2\,x+\sqrt{8Sy}.
\]
If $S=0$ then $\Q(B)=0$. Otherwise put $u=y/S\in[0,1]$, so $x=(1-u)S$
and the right side equals $S[1+\sqrt2(1-u+2\sqrt u)]$. Since
$1-u+2\sqrt u\le2$ on $[0,1]$, it is at most $(1+2\sqrt2)S$. Taking the
infimum over $\sigma$, and using $\Q(B)\le\rho$, gives
$\Q(B)\le(1+2\sqrt2)\min\{\rho,\inf_\sigma S\}=(1+2\sqrt2)\mathcal K(B)$.
\end{proof}

\subsection{Consequences}
\begin{corollary}[Selection form for finite laws]\label{cor:selection-form}
For every finite law, rarity and budget in Theorem~\ref{thm:direct-main},
\[
 10^{-20}\,\mathcal K(B)\le R_{\mu,\rho}(B)
 \le30001(1+2\sqrt2)\,\mathcal K(B)<1.2\cdot10^5\,\mathcal K(B).
\]
\end{corollary}

\begin{proof}
Combine Theorem~\ref{thm:direct-main} with
Proposition~\ref{prop:quantile-selection}.
\end{proof}

Corollary~\ref{cor:selection-form} is the finite-law case of the
characterization stated in~\cite{gl}; here it follows from
Theorem~\ref{thm:direct-main}, whose proof does not invoke the theorem
of~\cite{gl}. Conversely, that theorem, if granted, gives the rate
$R\asymp\Q$ of Theorem~\ref{thm:direct-main} for finite laws through
Proposition~\ref{prop:quantile-selection}, without the separate argument
of this paper. Its explicit constants $10^{-17}$ and $33000$ then give
$10^{-17}(1+2\sqrt2)^{-1}\Q\le R\le33000\,\Q$: a better lower constant
than $10^{-20}$ and a slightly worse upper one than $30001$.

The two results should be read as one rate established by two
arguments, not as independent characterizations. Nor are the two
arguments independent evidence: both manuscripts were developed from
the same research notes. They share the observable upper construction
(the upper bound here is the companion's construction applied to the
fixed quantile split) and most appendix lemmas, and differ mainly in the
lower-bound reduction.

\begin{remark}[General laws, conditionally]\label{rem:general-laws}
Proposition~\ref{prop:quantile-selection} holds for every probability
law on $(0,1/2]$, and $\Q(B)$ is then still a closed expression: $\kappa$
occupies finitely many bands. If $c\,\mathcal K\le R\le C\,\mathcal K$
holds for every such law, as~\cite{gl} states with $c=10^{-17}$ and
$C=33000$, then
\[
 \frac{c}{1+2\sqrt2}\,\Q(B)\le R_{\mu,\rho}(B)\le C\,\Q(B)
\]
for every probability law on $(0,1/2]$. This conditional extension is
not a result of the present manuscript; it rests entirely on the
theorem of~\cite{gl}.
\end{remark}

The comparison was supplied by an external adversarial review (Astra
Pro) of revision~1 of this manuscript, at repository commit
\texttt{2984835}. Its proof used hard selections; the proof above
works directly with the selection over submeasures of~\cite{gl}.
